% answers_short.tex —— 2013 年·北京卷·文科 · 简短答案速查
% 由 tools/make_short_answers.py 自 answers.tex 的 \gksolhead/\gkanswer 提取，
% 并与 2013/2013_数学真题_北京文_答案要点.md 逐题比对后生成，勿手改（改源 MD 后重跑）。

\begin{gknotice}
\textbf{使用说明}：本节只给最终答案便于快速核对，完整推导见解析版。答案为编写者独立推导，\textbf{非官方标准答案}。
\end{gknotice}

\gkpart{选择题}
\begin{center}\begin{tabular}{ccccccccc}
\toprule
题号 & 1 & 2 & 3 & 4 & 5 & 6 & 7 & 8 \\
\midrule
答案 & B & D & C & A & B & C & C & B \\
\bottomrule
\end{tabular}\end{center}

\gkpart{填空题}
\begin{itemize}[leftmargin=2.6em]
  \item[\textbf{9.}] $p=2$，准线方程为 $x=-1$。由 $\dfrac p2=1$ 得 $p=2$，准线 $x=-\dfrac p2$。
  \item[\textbf{10.}] $3$。底面是边长 $3$ 的正方形，高为 $1$，故 $V=\dfrac13\times3\times3\times1=3$。
  \item[\textbf{11.}] $q=2$，$S_n=2^{n+1}-2$。由 $a_3+a_5=q(a_2+a_4)$ 得 $q=2$，再由 $a_2+a_4=10a_1=20$ 得 $a_1=2$。
  \item[\textbf{12.}] $\dfrac{2\sqrt5}5$。点 $(1,0)$ 到边界 $2x-y=0$ 的距离为 $\dfrac2{\sqrt5}$，是到区域 $D$ 的最短距离。
  \item[\textbf{13.}] $(-\infty,2)$。$x\geqslant1$ 时 $\log_{\frac12}x\leqslant0$，$x<1$ 时 $2^x\in(0,2)$，两者合并即得。
  \item[\textbf{14.}] $3$。$D$ 是以 $\overrightarrow{AB}$ 与 $\overrightarrow{AC}$ 为邻边的平行四边形，面积 $\left|2\times2-1\times1\right|=3$。
\end{itemize}

\gkpart{解答题}
\begin{itemize}[leftmargin=2.6em]
  \item[\textbf{15.}] (1) 最小正周期 $T=\dfrac{\pi}2$，最大值 $\dfrac{\sqrt2}2$。$f(x)=\dfrac{\sqrt2}2\sin\left(4x+\dfrac{\pi}4\right)$；(2) $\dfrac{9\pi}{16}$。由 $\sin\left(4\alpha+\dfrac{\pi}4\right)=1$ 得 $\alpha=\dfrac{\pi}{16}+\dfrac{k\pi}2$，在 $\left(\dfrac{\pi}2,\pi\right)$ 内只有 $k=1$
  \item[\textbf{16.}] (1) $\dfrac6{13}$。$1$ 日至 $13$ 日中空气质量指数小于 $100$ 的有 $1$，$2$，$3$，$7$，$12$，$13$ 日共 $6$ 天；(2) $\dfrac4{13}$。两天中恰有 $1$ 天重度污染的到达日为 $4$，$5$，$7$，$8$ 日，共 $4$ 种；(3) 从 $5$ 日开始。$220$，$160$，$40$ 的方差为 $5600$，大于其余各起始日的三天方差
  \item[\textbf{17.}] (1) 证明见解析。由面面垂直的性质，$PA\subset$ 平面 $PAD$ 且 $PA\perp$ 交线 $AD$，故 $PA\perp$ 底面 $ABCD$；(2) 证明见解析。由 $AB\parallel DE$ 且 $AB=DE$ 得 $ABED$ 为平行四边形，故 $BE\parallel AD$；又 $AD\subset$ 平面 $PAD$，$BE\not\subset$ 平面 $PAD$，所以 $BE\parallel$ 平面 $PAD$；(3) 证明见解析。$BE\parallel AD$，$EF\parallel PD$，故平面 $BEF\parallel$ 平面 $PAD$，而 $CD\perp$ 平面 $PAD$ 给出平面 $PCD\perp$ 平面 $PAD$
  \item[\textbf{18.}] (1) $a=0$，$b=1$。$f'(x)=x(2+\cos x)$，切线水平要求 $f'(a)=0$，得 $a=0$，$b=f(0)=1$；(2) $(1,+\infty)$。$f$ 在 $x=0$ 处取最小值 $1$，且向两侧无界增大，故 $b>1$ 时恰有两个交点
  \item[\textbf{19.}] (1) $2\sqrt3$。菱形对角线互相垂直平分，$AC$ 为 $y=\dfrac12$ 被椭圆截得的弦；(2) 证明见解析。若为菱形则由 $|OA|=|AB|$、$|OC|=|CB|$ 与 $A$，$C$ 在 $W$ 上推出 $x_0y_0=0$，与 $B$ 不是顶点矛盾
  \item[\textbf{20.}] (1) $d_1=2$，$d_2=3$，$d_3=6$。$A_1=3$，$A_2=4$，$A_3=7$，而 $B_1=B_2=B_3=1$；(2) 证明见解析。数列递增，$A_i=a_i$，$B_i=a_{i+1}$，故 $d_i=a_1(1-q)q^{i-1}$ 是公比为 $q$ 的等比数列；(3) 证明见解析。由 $d_i$ 递增得 $a_2<\cdots<a_{n-1}$ 且 $a_n$ 为最小项，从而 $B_i=a_n$，$d_i=a_i-a_n$ 为等差数列
\end{itemize}
